% Beamer aligned with main.tex
% Compile from 3diags/: pdflatex SCRS_2026_3diagnostics_beamer.tex
% Figures: figs/ (same as main.tex; self-contained for Overleaf)

\PassOptionsToPackage{unicode}{hyperref}
\PassOptionsToPackage{hyphens}{url}
\documentclass[ignorenonframetext]{beamer}
\newif\ifbibliography

\setbeamertemplate{caption}[numbered]
\setbeamertemplate{caption label separator}{: }
\setbeamercolor{caption name}{fg=normal text.fg}
\beamertemplatenavigationsymbolsempty
\setbeamertemplate{part page}{%
  \centering
  \begin{beamercolorbox}[sep=16pt,center]{part title}
    \usebeamerfont{part title}\insertpart\par
  \end{beamercolorbox}}
\setbeamertemplate{section page}{%
  \centering
  \begin{beamercolorbox}[sep=12pt,center]{section title}
    \usebeamerfont{section title}\insertsection\par
  \end{beamercolorbox}}
\setbeamertemplate{subsection page}{%
  \centering
  \begin{beamercolorbox}[sep=8pt,center]{subsection title}
    \usebeamerfont{subsection title}\insertsubsection\par
  \end{beamercolorbox}}
\widowpenalties 1 10000
\raggedbottom
\AtBeginPart{\frame{\partpage}}
\AtBeginSection{%
  \ifbibliography\else\frame{\sectionpage}\fi}
\AtBeginSubsection{\frame{\subsectionpage}}

\usepackage{iftex}
\ifPDFTeX
  \usepackage[T1]{fontenc}
  \usepackage[utf8]{inputenc}
  \usepackage{textcomp}
\else
  \usepackage{unicode-math}
  \defaultfontfeatures{Scale=MatchLowercase}
  \defaultfontfeatures[\rmfamily]{Ligatures=TeX,Scale=1}
\fi
\usepackage{lmodern}
\usetheme{Madrid}
\usefonttheme{professionalfonts}
\IfFileExists{upquote.sty}{\usepackage{upquote}}{}
\IfFileExists{microtype.sty}{%
  \usepackage{microtype}
  \UseMicrotypeSet[protrusion]{basicmath}}{}
\makeatletter
\@ifundefined{KOMAClassName}{%
  \IfFileExists{parskip.sty}{\usepackage{parskip}}{%
    \setlength{\parindent}{0pt}%
    \setlength{\parskip}{6pt plus 2pt minus 1pt}}}{%
  \KOMAoptions{parskip=half}}
\makeatother
\setlength{\emergencystretch}{3em}
\providecommand{\tightlist}{%
  \setlength{\itemsep}{0pt}\setlength{\parskip}{0pt}}
\usepackage{booktabs}
\usepackage{graphicx}
\usepackage{appendixnumberbeamer}
\graphicspath{{figs/}}

\definecolor{ICLBlue}{HTML}{003E74}
\definecolor{ICLLightBlue}{HTML}{D4EFFC}
\definecolor{ICLTeal}{HTML}{02893B}
\setbeamercolor{palette primary}{bg=ICLBlue, fg=white}
\setbeamercolor{palette secondary}{bg=ICLBlue, fg=white}
\setbeamercolor{palette tertiary}{bg=ICLBlue, fg=white}
\setbeamercolor{titlelike}{fg=ICLBlue}
\setbeamercolor{title}{fg=white,bg=ICLBlue}
\setbeamercolor{subtitle}{fg=ICLBlue}
\setbeamercolor{author}{fg=ICLBlue}
\setbeamercolor{institute}{fg=ICLBlue}
\setbeamercolor{date}{fg=ICLBlue}
\setbeamercolor{frametitle}{fg=ICLBlue,bg=ICLLightBlue}
\setbeamerfont{frametitle}{series=\bfseries}
\setbeamercolor{structure}{fg=ICLBlue}
\setbeamercolor{section title}{fg=ICLBlue,bg=ICLLightBlue}
\setbeamercolor{subsection title}{fg=ICLBlue,bg=ICLLightBlue}
\setbeamercolor{item projected}{bg=ICLBlue, fg=white}
\setbeamertemplate{navigation symbols}{}

\usepackage{bookmark}
\IfFileExists{xurl.sty}{\usepackage{xurl}}{}
\urlstyle{same}
\hypersetup{%
  pdftitle={From Fitted to Fit-for-Purpose: three diagnostic strands},
  pdfauthor={Laurence T. Kell and Massimiliano Cardinale},
  hidelinks,
  pdfcreator={LaTeX via pandoc}}

\newcommand{\pfig}[2][]{%
  \includegraphics[#1]{#2}}

\title{\textbf{From Fitted to Fit-for-Purpose}}
\subtitle{Intertwining Internal, External, and Robustness Diagnostics\\[0.4em]
\small North Atlantic shortfin mako SS3 case study (SCRS/2026/074)}
\author{Laurence T. Kell \& Massimiliano Cardinale}
\institute{Imperial College London \and SLU Aqua, Lysekil}
\date{\today}

\begin{document}

\begin{frame}[plain]
  \titlepage
\end{frame}

% ---------------------------------------------------------------------------
\section{Framework}

\begin{frame}{Summary}
  \small
  Integrated assessments need more than convergence and fit.
  \begin{itemize}
    \item \textbf{Internal} --- consistency with likelihood, assumptions, and fitted data
    \item \textbf{External} --- consistency with specified evidence \emph{outside} the likelihood
    \item \textbf{Robustness} --- stability of advice under pre-specified alternative structures
  \end{itemize}
  \vspace{0.4em}
  Main text: framework + selected SS3 shortfin mako examples.\\
  Appendices: reproducible templates (plots, tables, \texttt{R} outputs).
  \vspace{0.4em}
  \begin{alertblock}{Best practice}
    Minimum diagnostic checklist; separate verification, validation, uncertainty, and sensitivity; rebuilding-based diagnostics for non-retained bycatch stocks.
  \end{alertblock}
\end{frame}

\begin{frame}{Why diagnostics beyond convergence?}
  \begin{itemize}
    \item Advice depends on latent quantities (SSB, recruitment, $F$) and equilibrium reference points
    \item SS3 flexibility increases the gap between a \emph{well-fitted} model and \emph{fit-for-purpose} advice
    \item A good fit alone is insufficient --- models should predict withheld data and be robust to structure
  \end{itemize}
  \vspace{0.4em}
  \begin{alertblock}{Three intertwined strands}
    Internal \textbullet\ External \textbullet\ Robustness --- complementary, not hierarchical
  \end{alertblock}
\end{frame}

\begin{frame}{Three diagnostic strands}
  \small
  \begin{table}
    \centering
    \begin{tabular}{@{}p{0.20\linewidth}p{0.36\linewidth}p{0.38\linewidth}@{}}
      \toprule
      \textbf{Class} & \textbf{Main question} & \textbf{SS3 examples} \\
      \midrule
      Internal &
      Consistent with likelihood and fitted data? &
      Convergence, Hessian, residuals, ESS, rec.\ devs, profiles \\
      External &
      Consistent with evidence outside likelihood? &
      Life history, maturity--selectivity, survey perception \\
      Robustness &
      Advice stable under pre-specified alternatives? &
      Retros, hindcasts, ensembles, sensitivities \\
      \bottomrule
    \end{tabular}
  \end{table}
  \vspace{0.4em}
  \begin{itemize}
    \item \textbf{Verification} --- implemented correctly and reproducibly?
    \item \textbf{Validation} --- adequate for management use (hindcasts)?
    \item \textbf{Uncertainty} (within model) vs.\ \textbf{sensitivity} (across structures)
  \end{itemize}
\end{frame}

% ---------------------------------------------------------------------------
\section{Internal diagnostics}

\begin{frame}{Internal diagnostics}
  \textbf{Beyond convergence}
  \begin{itemize}
    \item What data inform steepness, $M$, selectivity, catchability --- estimated, fixed, or prior-driven?
    \item Distinguish \textbf{identifiability}, \textbf{estimability}, and \textbf{regularisation}
    \item Likelihood profiles and Hessian checks; MCMC/\texttt{adnuts} when weakly identified
  \end{itemize}
  \vspace{0.4em}
  \textbf{Minimum checklist}
  \begin{itemize}
    \item Convergence, gradient, Hessian, parameters near bounds
    \item Residuals by data source; composition ESS; recruitment-deviation structure
  \end{itemize}
\end{frame}

\begin{frame}{Process error diagnostics}
  \small
  Do production-function residuals behave like white noise, or reveal misspecification / non-stationarity?
  \begin{enumerate}
    \item \textbf{Temporal patterns} --- trends, regime shifts (STARS)
    \item \textbf{Biomass-dependent bias} --- shape misspecification, heteroscedasticity
    \item \textbf{ACF} --- white vs.\ \textbf{red noise}; uncertainty intervals may be too narrow
    \item \textbf{Distribution} --- skewness, heavy tails, offset mean
  \end{enumerate}
  \vspace{0.3em}
  In SS3, misspecification in $M$, selectivity, catchability is often absorbed into \textbf{recruitment deviations}.
\end{frame}

\begin{frame}{Internal vs.\ external stock--recruitment fitting}
  \small
  \textbf{Internal} (e.g.\ SAM): S--R estimated in state space with SSB--recruitment uncertainty correlated.
  \vspace{0.3em}
  \textbf{External post-hoc} (common with SS3 outputs):
  \begin{itemize}
    \item SSB and recruitment are model-dependent point estimates
    \item Ignores estimation uncertainty and shared errors
    \item Apparent S--R patterns can reflect assessment structure, not biology
    \item ``Quirky'' recruitment series may be assessment artefacts
  \end{itemize}
  \vspace{0.3em}
  Process error from external S--R fits may reflect \emph{model structure}, not true stochasticity.
\end{frame}

\begin{frame}{Production functions}
  \[
    SP_t = B_{t+1} - B_t + C_t
  \]
  \textbf{Explicit} (JABBA, SPiCT): production function estimated directly\\
  \textbf{Implicit} (SS3): emergent from growth, recruitment, $M$, selectivity
  \vspace{0.5em}
  \begin{block}{Diagnostic classes are model-agnostic}
    Implementation differs; the three strands apply to any assessment framework.
  \end{block}
\end{frame}

% ---------------------------------------------------------------------------
\section{Case study}

\begin{frame}{Shortfin mako SS3 scenarios}
  \small
  Illustrative scenarios from the 2025 ICCAT SMA meeting; \textbf{not adopted for advice}.
  \begin{itemize}
    \item Catches dominated by juveniles/subadults; management state = mature-female SSB
    \item Stress test: separate internal fit, external biological consistency, robustness of interpretation
    \item Workflows: \texttt{r4ss}, \texttt{ss3diags}, \texttt{sstruth}, \texttt{FLRebuild}, \texttt{kobe}
  \end{itemize}
\end{frame}

\begin{frame}{Internal case study: MSY-scaled trajectories}
  \centering
  \pfig[width=0.92\linewidth,height=0.78\textheight,keepaspectratio]{sma-fig-ts-1.png}
  \vspace{-0.2em}
  {\scriptsize Catch tracks $F$ more closely than SSB; implied dynamics vs.\ reference points need scrutiny.}
\end{frame}

\begin{frame}{Internal case study: production phase plot}
  \centering
  \pfig[width=0.92\linewidth,height=0.78\textheight,keepaspectratio]{sma-fig-pe-phase-1.png}
  \vspace{-0.2em}
  {\scriptsize Realised surplus production does not always align with equilibrium curves.}
\end{frame}

\begin{frame}{External case study: maturity vs.\ selectivity}
  \centering
  \pfig[width=0.88\linewidth,height=0.72\textheight,keepaspectratio]{selmat.png}
  \vspace{0.3em}
  \begin{block}{Currency mismatch}
    $F$ on immature ages; state on mature-female SSB --- weak surplus production on SSB may reflect age structure, not biology alone.
  \end{block}
\end{frame}

% ---------------------------------------------------------------------------
\section{Robustness}

\begin{frame}{Robustness diagnostics}
  \textbf{Separate}
  \begin{itemize}
    \item Estimation uncertainty \emph{conditional on} one accepted run
    \item Structural sensitivity \emph{across} pre-specified scenarios
  \end{itemize}
  \vspace{0.3em}
  \textbf{Report both}
  \begin{itemize}
    \item \textbf{Central tendencies} for targets ($B/B_{MSY}$, $F/F_{MSY}$, catch)
    \item \textbf{Tail probabilities} for limits ($P(B<B_{lim})$, $P(F>F_{lim})$)
  \end{itemize}
  \vspace{0.3em}
  \textbf{Cross-framework}: SPiCT/JABBA (explicit process error) vs.\ SS3/SAM (recruitment-dev partition) --- compare on common SSB surplus-production currency in ensembles.
\end{frame}

\begin{frame}{Robustness case study: residuals over time}
  \centering
  \pfig[width=0.92\linewidth,height=0.72\textheight,keepaspectratio]{sma-fig-res-time-1.png}
  \vspace{-0.2em}
  {\scriptsize STARS-type shading: candidate regime shifts in process residuals.}
\end{frame}

\begin{frame}{Robustness case study: autocorrelation}
  \centering
  \pfig[width=0.92\linewidth,height=0.62\textheight,keepaspectratio]{sma-fig-res-acf-1.png}
  \vspace{0.3em}
  \small
  Low-lag ACF outside 95\% bounds $\Rightarrow$ \textbf{red noise}.\\
  Forecast and rebuilding intervals based on white-noise assumptions may \textbf{underestimate risk}.
\end{frame}

\begin{frame}[shrink=5]{Red noise and heteroscedasticity}
  \small
  \textbf{Red noise} --- environmental forcing, omitted covariates, serial catch/effort errors, age-structure effects.
  \vspace{0.3em}
  \textbf{Heteroscedasticity} --- clustered large residuals in time or at low biomass:
  \begin{itemize}
    \item Misreporting, selectivity shifts treated as fixed
    \item For bycatch stocks: higher variance at low biomass may be genuine demographic stochasticity
  \end{itemize}
  \vspace{0.3em}
  \begin{alertblock}{Message}
    Equilibrium MSY reference points from a single SS3 run are \textbf{conditional} --- gate with PE diagnostics, hindcasts, and ensembles.
  \end{alertblock}
\end{frame}

\begin{frame}{Four-panel PE diagnostics}
  \centering
  \pfig[width=0.92\linewidth,height=0.62\textheight,keepaspectratio]{4panel.png}
  \vspace{0.3em}
  {\small Time \textbullet\ Biomass \textbullet\ ACF \textbullet\ Histogram --- standard internal/robustness panel set.}
\end{frame}

\begin{frame}{Retrospectives then hindcasts}
  \textbf{Two-step workflow} (step~2 is forecast-only, \texttt{-nohess}, parallel peels):
  \begin{enumerate}
    \item \textbf{Retrospectives} --- full SS3 re-fit per peel
    \item \textbf{Catch back-tests} --- forecast from peel \texttt{ss3.par}\\
          \texttt{catchReported}, \texttt{fMeanFull}, \texttt{fZero}, \texttt{fMsy}
  \end{enumerate}
  \vspace{0.4em}
  \small
  One-step prediction errors from retros; management-scenario forecasts from hindcasts.\\
  Status: 88/88 forecast runs collected (Scenarios 2--3).
\end{frame}

\begin{frame}[shrink=8]{Ensemble weights (Scenarios 1--3)}
  \small
  \begin{table}
    \centering
    \begin{tabular}{lrrrr}
      \toprule
      Scenario & Equal & ICES (Akaike) & CVindexmase & CVindexrho \\
      \midrule
      Scenario-1 & 0.33 & $\approx 0$ & --- & --- \\
      Scenario-2 & 0.33 & $\approx 1$ & 1.00 & 1.00 \\
      Scenario-3 & 0.33 & $\approx 0$ & $\approx 0$ & 0 \\
      \bottomrule
    \end{tabular}
  \end{table}
  \vspace{0.3em}
  MVLN ensemble combines scenario draws; compare equal, likelihood, and hindcast-skill weights.
\end{frame}

% ---------------------------------------------------------------------------
\section{Lessons and discussion}

\begin{frame}[shrink=8]{SCRS diagnostic checklist}
  \small
  \begin{enumerate}
    \item \textbf{Internal:} convergence, Hessian, bounds, residuals, ESS, information content
    \item \textbf{External:} maturity--selectivity, survey perception, signals outside likelihood
    \item \textbf{Robustness:} retros, hindcasts, sensitivities, ensembles; tails (limits) and medians (targets)
    \item \textbf{Process error:} STARS/regimes, biomass bias, ACF, distribution; compare SS3 vs.\ JABBA/SPiCT partitions
  \end{enumerate}
\end{frame}

\begin{frame}{Rebuilding-based diagnostics (bycatch)}
  \small
  For non-retained shortfin mako, equilibrium $B_{MSY}$ / $F_{MSY}$ may be a poor primary objective.
  \vspace{0.3em}
  \begin{itemize}
    \item Juvenile-dominated catches vs.\ mature-female SSB state
    \item Report $B_{\mathrm{rebuild}}$, $T_{\mathrm{rebuild}}$, $P(B_T \ge B_{\mathrm{target}})$ under stated bycatch scenarios
    \item Three strands map to production behaviour (internal), age structure (external), projection stability (robustness)
  \end{itemize}
\end{frame}

\begin{frame}{Observation vs.\ process error}
  \small
  \textbf{Observation error} --- survey/measurement variation.\\
  \textbf{Process error} --- environmental and demographic stochasticity, omitted biology.
  \vspace{0.3em}
  \begin{itemize}
    \item SS3: input CVs for observations; most process variability in rec.\ devs
    \item SPiCT: joint estimation of process and observation variance
  \end{itemize}
  \vspace{0.3em}
  Practical checks: index residual structure; consistency across data types; hindcast skill.\\
  If partition is uncertain $\Rightarrow$ conservative uncertainty and cross-framework ensembles.
\end{frame}

\begin{frame}{Reproducible appendices (\texttt{3diags/})}
  \small
  \begin{itemize}
    \item \texttt{main.tex} + three appendix strands (\texttt{render\_appendices.R})
    \item Internal / external / robustness templates with case-study outputs
    \item Long runs: retros \checkmark; catch back-tests \checkmark; ensemble weights \checkmark; LOO CV setup only
  \end{itemize}
  \vspace{0.4em}
  Packages should hold generic logic; appendix Rmd files remain thin case-study wrappers.
\end{frame}

\begin{frame}{Validation and advice implications}
  \small
  \begin{itemize}
    \item Validate against observations where possible --- not only model estimates
    \item Joint CPUE is a structural choice; alternative constructions as ensemble members
    \item SS version can shift MSY quantities despite similar likelihoods
    \item No single SS3 run to advice without agreed diagnostics
    \item Non-stationary PE $\Rightarrow$ alternative hypotheses and ensembles
  \end{itemize}
\end{frame}

\begin{frame}{Conclusions}
  \begin{itemize}
    \item Internal, external, and robustness diagnostics answer different questions
    \item Shortfin mako illustrates age-structure and currency effects on PE interpretation
    \item Routine SCRS reporting: three strands together; verification / validation / uncertainty / sensitivity separated
    \item Reproducible appendices support WG review without overloading the main paper
  \end{itemize}
\end{frame}

\begin{frame}{Acknowledgements}
  \small
  We thank ICCAT and participants in the 2025 Shortfin Mako Shark Stock Assessment Meeting,
  and the 2026 Working Group on Stock Assessment Methods.
  LK acknowledges support from Spanish surface longline producer organisations OPP-Burela,
  OPNAPA, OPROMAR, and ORPAGU under the 2026 collaboration agreement.
\end{frame}

\end{document}
